% Nouveau EGA V (cote 25): archive cote 25, site page 4 (PDF page 5).
% typescript. Draft compared: none.
% AI-transcribed by Claude Opus 5.5 (claude-opus-5-5, adaptive thinking, effort high), Message Batches API, batch pilot (msgbatch_0189Nhf1xB615YWnztubtLfc), 2026-11-12; not yet checked by eye.
\scanpage{25/4}
IV 20 \quad \uline{Sections hyperplanes et projections \deleted{linéaires}\inserted{coniques}}.

20.1. \quad \uline{Préliminaires \over{de}{\unsure{et}} notations}.

Soient S un \quad préschéma , $\underline{E}$ un Module loc libre de type fini sur S,
$\check{\underline{E}}$ son dual. Nous noterons par P=$\underline{P}(\underline{E})$ le fibré projectif défini par $\underline{E}$,
par $\check{P}$ celui défini par $\check{\underline{E}}$ . Nous\hand{|}appelerons \ $\check{P}$ \ le \uline{schéma des hyperplans}
de P, cette terminologie se justifiant de la façon suivante. \ Soit
$\hand{\xi}$ une section de $\check{P}$ sur S, qui est donc\hand{/}déterminée par un module quotient
inversible L de $\check{\underline{E}}$ . On en conclut un module quotient inversible
$\underline{L}_P$ de$(\check{\underline{E}})_P$ = $(\underline{E}_P)^{\hand{\vee}}$ , d'autre part on dispose d'un module quotient inversibl
$\underline{O}_P(1)$ de \ $\underline{E}_P$ . Passant \ aux duals, on peut envisager \ $\underline{L}_P^{-1}$ resp. $\underline{O}_P(-1)$
comme des sous-modules inversibles (loc facteurs directs) de \ $\underline{E}_P$ resp.
de $(\underline{E}_P)^{\hand{\vee}}$ , \ et l'accouplement \ $\underline{E}_P \otimes \check{\underline{E}}_P \longrightarrow \underline{O}_P$ \ définit \ donc un
accouplement naturel
(x) \[ \underline{O}_P(-1)\otimes\underline{L}_P^{-1} \to \underline{O}_P \]
ou encore un homomorphisme \ transposé
(xx) \[ \underline{O}_P \longrightarrow \underline{O}_P(1)\otimes\underline{L}_P \ = \ \unsure{\underline{L}_P(1)|\underline{D}_P(1)} \ , \]
i.e. une section de $\underline{L}_P(1)$, canoniquement déterminée par $\hand{\xi}$ . Le''diviseur''
de cette section, \ i.e. le sous-schéma\inserted{fermé} $H_{\hand{\xi}}$ de \ P défini par l'Idéal
image de (x), \ est \ appelé l'\uline{hyperplan dans P} \ défini par \ l'élément
$\hand{\xi} \hand{\in} \check{P}(S)$ . \ Il peut encore se décrire en notant que localement sur S \xout{la}
\xout{\unsure{donnée de $\check{\xi}$ équivaut}} \quad $\hand{\xi}$ \ est donné par une section \ $\phi$ de \ $\underline{E}$ , \xout{de} telle
\xout{\illegible{1}} que $\phi(s) \neq 0$ pour tout s \ ( $\phi$ étant déterminé par $\hand{\xi}$ \ modulo
multiplication par une section inversible de $\underline{O}_S$ ); comme $\underline{E}$ = $p_{\hand{*}}(\underline{O}_P(1))$
( p: P$\to$S étant la projection), \ $\phi$ \ peut être considéré comme une section
de \ $\underline{O}_P(1)$, don\over{c}{t} le diviseur n'est autre que $H_{\hand{\xi}}$. \ Bien entendu, si on
\inserted[typed]{\uline{considére}}\xout{\illegible{1}} $\underline{L}^{-1}$ \ comme un sous-\xout{\illegible{1}} module inversible de \ $\underline{E}$ locale
ment facteur direct dans $\underline{E}$, alors la correspondance entre $\hand{\xi}$ \ (i.e. $\underline{L}$ \hand{\unsure{$\xi$,}}

\mspage{\hand{4}}
